\scanpage{177}
\deleted{\illegible{1} loc. sur $X^{(p/S)},\ \underline O_X$} Prenons $L=\underline O_X$, avec $\Delta_0$ la connexion habituelle. Alors $\omega$ satisfaisant $d\omega=0$, $\widetilde C\omega=\widetilde\omega$ définit donc une donnée de descente sur $L$, qui \deleted{qui} provient de $L_0$ sur $X_0^{p/S}=(X^{p/S},\underline\Lambda)$. Localement, on peut alors \deleted{\unsure{splitter cette connexion,}} trouver une base de $L_0$, i.e. loclt on peut engendrer $L$ par une section horizontale $f$, ce qui signifie que loc. on a $\omega=\frac{df}{f}$. On trouve ainsi le

\uline{Théorème}\quad Pour que \inserted{la 1-forme} $\omega$ soit loc. une dérivée logarithmique, $\omega=df/f$ ($f\in\Gamma(O_X^*)$) il f. et s. que l'on ait $d\omega=0$, \inserted{$C\omega=\widetilde\omega$.} En d'autres termes, on a un hom $f\longrightarrow df/f$ :
\[
 O_X^*\longrightarrow\underline\Omega^1_{X/S}
\]
dont l'image est le sous-faisceau de $\underline\Omega^1_{X/S}$ formé des $\omega$ telles que $d\omega=0$, $\widetilde C\omega=\widetilde\omega$. [NB le noyau est $\deleted{\unsure{\underline O}}\Lambda^*$, donc si $S$ est parfait, c'est $\underline O_X^{*\,p}$]

\textcircled{6.} \uline{Cas des schémas relatifs lisses. Relations avec la trace des formes diff. de degré maximum.}

Supposons $X$ lisse sur $S$, donc $X^{p/S}$ lisse sur $S$. Alors $X\xrightarrow{f_{X/S}}X^{p/S}$ est un morphisme fini localmt libre (en plus d'être radiciel et surjectif), à
% archive sheet number: 176

